\documentclass{beamer}
\usetheme{Madrid}
\usecolortheme{default}

% Define custom colors for AI slides
\definecolor{AIBlue}{RGB}{30, 144, 255}
\definecolor{AILightBlue}{RGB}{173, 216, 230}
\definecolor{AINavy}{RGB}{25, 25, 112}

% Commands to switch to AI color scheme
\newcommand{\AIcolors}{%
  \setbeamercolor{frametitle}{bg=AIBlue,fg=white}%
  \setbeamercolor{title}{bg=AINavy,fg=white}%
  \setbeamercolor{structure}{fg=AIBlue}%
  \setbeamercolor{item}{fg=AINavy}%
  \setbeamercolor{normal text}{fg=black}%
}

% Command to restore default colors
\newcommand{\defaultcolors}{%
  \setbeamercolor{frametitle}{bg=osaorange,fg=white}%
  \setbeamercolor{title}{bg=osaorange,fg=white}%
  \setbeamercolor{structure}{fg=osaorange}%
  \setbeamercolor{item}{fg=black}%
  \setbeamercolor{normal text}{fg=black}%
}

% Define OSU orange color
\definecolor{osaorange}{RGB}{220,115,38}

% Set colors
\setbeamercolor{structure}{fg=osaorange}
\setbeamercolor{title}{fg=white,bg=osaorange}
\setbeamercolor{frametitle}{fg=white,bg=osaorange}

\usepackage{graphicx}
\usepackage{amsmath}
\usepackage{booktabs}
\usepackage{hyperref}
\usepackage{tikz}
\usetikzlibrary{shapes,arrows}

\title{H524 Introduction to Biostatistics}
\subtitle{Week 2: Probability Concepts and Applications}
\author{Dr. John Molitor}
\institute{Oregon State University\\College of Health}
\date{Fall 2025}

\begin{document}

\frame{\titlepage}

\begin{frame}
\frametitle{Outline}
\tableofcontents
\end{frame}

\section{Review from Week 1}

\begin{frame}
\frametitle{Quick Review: Key Concepts}
\begin{itemize}
\item \textbf{Data Types:} Quantitative vs Qualitative
\item \textbf{Data Presentation:} Histograms, box plots, bar charts
\item \textbf{Summary Measures:} Central tendency and variability
\item \textbf{Rates and Standardization:} Age-specific rates, direct/indirect methods
\end{itemize}

\vspace{0.5cm}
\textbf{This Week:} Chapter 6 - Probability foundations for statistical inference

\vspace{0.3cm}
\textbf{Reading:} Pagano \& Gauvreau Chapter 6 (6.1-6.4.2, 6.5)
\end{frame}

\section{Basic Probability Concepts (Ch 6.1)}

\begin{frame}
\frametitle{Why Probability in Biostatistics?}
\begin{itemize}
\item \textbf{Foundation for Statistical Inference}
  \begin{itemize}
  \item Understanding uncertainty in medical research
  \item Bridge between sample data and population conclusions
  \end{itemize}
\item \textbf{Clinical Applications}
  \begin{itemize}
  \item Diagnostic test interpretation
  \item Treatment efficacy assessment
  \item Risk communication to patients
  \end{itemize}
\item \textbf{Research Design}
  \begin{itemize}
  \item Power calculations
  \item Sample size determination
  \item Study planning and interpretation
  \end{itemize}
\end{itemize}
\end{frame}

\begin{frame}
\frametitle{Basic Probability Concepts}
\begin{itemize}
\item \textbf{Experiment:} Any process that generates observations
  \begin{itemize}
  \item Example: Testing a patient for COVID-19
  \end{itemize}
\item \textbf{Sample Space ($S$):} Set of all possible outcomes
  \begin{itemize}
  \item Example: $S = \{\text{Positive}, \text{Negative}\}$
  \end{itemize}
\item \textbf{Event:} Subset of the sample space
  \begin{itemize}
  \item Example: Patient tests positive
  \end{itemize}
\item \textbf{Probability:} Measure of likelihood (0 to 1)
  \begin{itemize}
  \item $P(\text{Event}) = \frac{\text{Number of favorable outcomes}}{\text{Total number of possible outcomes}}$
  \end{itemize}
\end{itemize}
\end{frame}

\begin{frame}
\frametitle{Probability Rules}
\begin{enumerate}
\item \textbf{Range:} $0 \leq P(A) \leq 1$ for any event $A$
\item \textbf{Certainty:} $P(S) = 1$ (something must happen)
\item \textbf{Complement:} $P(A^c) = 1 - P(A)$
\item \textbf{Addition Rule:}
  $$P(A \cup B) = P(A) + P(B) - P(A \cap B)$$
\item \textbf{Multiplication Rule:}
  $$P(A \cap B) = P(A) \times P(B|A)$$
\end{enumerate}
\end{frame}

\begin{frame}
\frametitle{Medical Example: Diagnostic Testing}
\textbf{COVID-19 Rapid Test Performance:}
\begin{itemize}
\item Sensitivity = 85\% (probability of positive test given disease)
\item Specificity = 97\% (probability of negative test given no disease)
\item Disease prevalence = 5\%
\end{itemize}

\vspace{0.5cm}
\textbf{Questions:}
\begin{itemize}
\item What's the probability of a false positive?
\item What's the probability someone with a positive test has COVID?
\item How does prevalence affect test interpretation?
\end{itemize}

\vspace{0.5cm}
These questions require understanding conditional probability!
\end{frame}

\section{Conditional Probability (Ch 6.2)}

\begin{frame}
\frametitle{Conditional Probability}
\textbf{Definition:} Probability of event $A$ given that event $B$ has occurred

$$P(A|B) = \frac{P(A \cap B)}{P(B)}, \text{ where } P(B) > 0$$

\vspace{0.5cm}
\textbf{Medical Example:}
\begin{itemize}
\item $A$ = Patient has diabetes
\item $B$ = Patient is overweight
\item $P(A|B)$ = Probability of diabetes given patient is overweight
\end{itemize}

\vspace{0.5cm}
\textbf{Key Insight:} Conditional probability updates our beliefs based on new information
\end{frame}

\begin{frame}
\frametitle{Independence}
\textbf{Definition:} Events $A$ and $B$ are independent if:
$$P(A|B) = P(A) \text{ or equivalently } P(A \cap B) = P(A) \times P(B)$$

\vspace{0.5cm}
\textbf{Examples in Healthcare:}
\begin{itemize}
\item \textbf{Independent:}
  \begin{itemize}
  \item Blood type and eye color
  \item Two different patients' test results
  \end{itemize}
\item \textbf{Dependent:}
  \begin{itemize}
  \item Smoking status and lung cancer
  \item Age and cardiovascular disease
  \item Family history and genetic disorders
  \end{itemize}
\end{itemize}

\vspace{0.5cm}
\textbf{Caution:} Independence is often assumed but should be verified!
\end{frame}

\begin{frame}
\frametitle{Diagnostic Test Performance}
\textbf{Key Metrics for Medical Tests:}
\begin{itemize}
\item \textbf{Sensitivity:} P(Positive Test | Disease Present)
  \begin{itemize}
  \item Ability to correctly identify disease
  \item True positive rate
  \end{itemize}
\item \textbf{Specificity:} P(Negative Test | Disease Absent)
  \begin{itemize}
  \item Ability to correctly identify no disease
  \item True negative rate
  \end{itemize}
\end{itemize}

\vspace{0.5cm}
\textbf{Example - COVID-19 Rapid Test:}
\begin{itemize}
\item Sensitivity = 85\% (detects 85\% of infected people)
\item Specificity = 97\% (correctly identifies 97\% of non-infected)
\end{itemize}
\end{frame}

\begin{frame}
\frametitle{Predictive Values}
\textbf{Clinical Interpretation Measures:}
\begin{itemize}
\item \textbf{Positive Predictive Value (PPV):} P(Disease | Positive Test)
  \begin{itemize}
  \item Probability patient has disease given positive test
  \item Depends on disease prevalence
  \end{itemize}
\item \textbf{Negative Predictive Value (NPV):} P(No Disease | Negative Test)
  \begin{itemize}
  \item Probability patient is healthy given negative test
  \item Also depends on disease prevalence
  \end{itemize}
\end{itemize}

\vspace{0.5cm}
\textbf{Key Insight:} PPV and NPV change with disease prevalence in the population, while sensitivity and specificity are test characteristics.
\end{frame}

\section{Diagnostic Tests (Ch 6.3)}

\begin{frame}
\frametitle{2x2 Contingency Tables}
\begin{center}
\begin{tabular}{|c|c|c|c|}
\hline
& \textbf{Disease +} & \textbf{Disease -} & \textbf{Total} \\
\hline
\textbf{Exposed} & a & b & a + b \\
\hline
\textbf{Not Exposed} & c & d & c + d \\
\hline
\textbf{Total} & a + c & b + d & n \\
\hline
\end{tabular}
\end{center}

\vspace{0.5cm}
\textbf{Key Measures:}
\begin{itemize}
\item \textbf{Risk in Exposed:} $R_E = \frac{a}{a+b}$
\item \textbf{Risk in Unexposed:} $R_U = \frac{c}{c+d}$
\item \textbf{Relative Risk (RR):} $\frac{R_E}{R_U} = \frac{a/(a+b)}{c/(c+d)}$
\item \textbf{Odds Ratio (OR):} $\frac{ad}{bc}$
\end{itemize}
\end{frame}

\begin{frame}
\frametitle{Relative Risk vs Odds Ratio}
\begin{itemize}
\item \textbf{Relative Risk (RR):}
  \begin{itemize}
  \item Direct measure of risk comparison
  \item Used in cohort studies
  \item RR = 2 means exposed have twice the risk
  \item Cannot be calculated in case-control studies
  \end{itemize}
\item \textbf{Odds Ratio (OR):}
  \begin{itemize}
  \item Approximates RR when disease is rare
  \item Can be calculated in all study designs
  \item More difficult to interpret directly
  \item OR > RR when outcome is common
  \end{itemize}
\end{itemize}

\vspace{0.5cm}
\textbf{Rule of thumb:} When disease prevalence < 10\%, OR $\approx$ RR
\end{frame}

\begin{frame}
\frametitle{Example: Smoking and Lung Cancer}
\textbf{Study of 1000 individuals over 20 years:}

\begin{center}
\begin{tabular}{|l|c|c|c|}
\hline
& \textbf{Lung Cancer} & \textbf{No Cancer} & \textbf{Total} \\
\hline
\textbf{Smokers} & 60 & 440 & 500 \\
\hline
\textbf{Non-smokers} & 10 & 490 & 500 \\
\hline
\textbf{Total} & 70 & 930 & 1000 \\
\hline
\end{tabular}
\end{center}

\vspace{0.3cm}
\textbf{Calculations:}
\begin{itemize}
\item Risk in smokers: $60/500 = 0.12$ (12\%)
\item Risk in non-smokers: $10/500 = 0.02$ (2\%)
\item Relative Risk: $0.12/0.02 = 6.0$
\item Odds Ratio: $(60 \times 490)/(10 \times 440) = 6.68$
\end{itemize}

\textbf{Interpretation:} Smokers have 6 times the risk of lung cancer
\end{frame}

\section{Random Variables and Expected Value (Ch 6.4)}

\begin{frame}
\frametitle{Expected Value}
\textbf{Definition:} Long-run average value of a random variable

\vspace{0.5cm}
\textbf{Discrete Random Variable:}
$$E(X) = \mu = \sum_{i} x_i \cdot P(X = x_i)$$

\vspace{0.5cm}
\textbf{Medical Example:} Number of adverse events in clinical trial
\begin{center}
\begin{tabular}{|c|c|c|c|c|}
\hline
Events ($x$) & 0 & 1 & 2 & 3 \\
\hline
Probability & 0.70 & 0.20 & 0.08 & 0.02 \\
\hline
\end{tabular}
\end{center}

$$E(X) = 0(0.70) + 1(0.20) + 2(0.08) + 3(0.02) = 0.42$$

\textbf{Interpretation:} On average, expect 0.42 adverse events per patient
\end{frame}

\begin{frame}
\frametitle{Variance and Standard Deviation}
\textbf{Variance:} Measure of spread around the expected value
$$\text{Var}(X) = \sigma^2 = E[(X - \mu)^2] = E(X^2) - [E(X)]^2$$

\textbf{Standard Deviation:} Square root of variance
$$SD(X) = \sigma = \sqrt{\text{Var}(X)}$$

\vspace{0.5cm}
\textbf{Properties:}
\begin{itemize}
\item $\text{Var}(aX) = a^2 \text{Var}(X)$
\item $\text{Var}(X + Y) = \text{Var}(X) + \text{Var}(Y)$ if $X$ and $Y$ are independent
\item Standard deviation has same units as original data
\end{itemize}
\end{frame}

\section{Discrete Distributions (Ch 6.5)}

\begin{frame}
\frametitle{Binomial Distribution}
\textbf{Conditions:}
\begin{itemize}
\item Fixed number of trials ($n$)
\item Two possible outcomes (success/failure)
\item Constant probability of success ($p$)
\item Independent trials
\end{itemize}

\vspace{0.5cm}
\textbf{Probability Mass Function:}
$$P(X = k) = \binom{n}{k} p^k (1-p)^{n-k}$$

\vspace{0.5cm}
\textbf{Parameters:}
\begin{itemize}
\item Mean: $\mu = np$
\item Variance: $\sigma^2 = np(1-p)$
\end{itemize}
\end{frame}

\begin{frame}
\frametitle{Binomial Example: Vaccine Efficacy}
\textbf{Clinical Trial:} 20 vaccinated individuals exposed to virus
\begin{itemize}
\item Vaccine efficacy = 90\% (probability of protection)
\item What's the probability exactly 18 are protected?
\end{itemize}

\vspace{0.5cm}
\textbf{Solution:}
\begin{itemize}
\item $n = 20$, $p = 0.90$, $k = 18$
\item $P(X = 18) = \binom{20}{18} (0.90)^{18} (0.10)^2$
\item $P(X = 18) = 190 \times 0.1501 \times 0.01 = 0.285$
\end{itemize}

\vspace{0.5cm}
\textbf{Expected number protected:} $\mu = 20 \times 0.90 = 18$\\
\textbf{Standard deviation:} $\sigma = \sqrt{20 \times 0.90 \times 0.10} = 1.34$
\end{frame}

\begin{frame}
\frametitle{Normal Distribution}
\textbf{Properties:}
\begin{itemize}
\item Bell-shaped, symmetric curve
\item Characterized by mean ($\mu$) and standard deviation ($\sigma$)
\item Area under curve = 1
\item Approximately 68\% within $\pm 1\sigma$
\item Approximately 95\% within $\pm 2\sigma$
\item Approximately 99.7\% within $\pm 3\sigma$
\end{itemize}

\vspace{0.5cm}
\textbf{Standard Normal:} $Z = \frac{X - \mu}{\sigma}$ has mean 0, SD 1

\vspace{0.5cm}
\textbf{Medical Applications:}
\begin{itemize}
\item Blood pressure measurements
\item Cholesterol levels
\item Height and weight
\item Many lab values
\end{itemize}
\end{frame}

\begin{frame}
\frametitle{Normal Distribution Example}
\textbf{Systolic Blood Pressure in Adults:}
\begin{itemize}
\item Mean = 120 mmHg
\item Standard deviation = 15 mmHg
\item What proportion have BP > 140 mmHg (hypertension)?
\end{itemize}

\vspace{0.5cm}
\textbf{Solution:}
\begin{enumerate}
\item Standardize: $Z = \frac{140 - 120}{15} = 1.33$
\item Look up in Z-table: $P(Z < 1.33) = 0.9082$
\item Therefore: $P(BP > 140) = 1 - 0.9082 = 0.0918$
\end{enumerate}

\vspace{0.5cm}
\textbf{Interpretation:} About 9.2\% of adults have hypertension by this criterion
\end{frame}

\section{Applications and Examples}

\begin{frame}
\frametitle{Probability in Public Health}
\textbf{Disease Surveillance Example:}
\begin{itemize}
\item Population: 50,000 individuals
\item New disease cases: 125 per year
\item Deaths from disease: 15 per year
\end{itemize}

\vspace{0.5cm}
\textbf{Calculate probabilities:}
\begin{itemize}
\item P(disease) = 125/50,000 = 0.0025 (0.25\%)
\item P(death | disease) = 15/125 = 0.12 (12\%)
\item P(death and disease) = 15/50,000 = 0.0003 (0.03\%)
\end{itemize}

\vspace{0.5cm}
\textbf{Clinical interpretation:}
\begin{itemize}
\item Incidence rate: 250 per 100,000 per year
\item Case fatality rate: 12\%
\item Overall mortality rate: 30 per 100,000 per year
\end{itemize}
\end{frame}

\section{AI as a Learning Tool for Probability}

\AIcolors
\begin{frame}
\frametitle{AI for Probability Calculations}
\textbf{How AI can help with probability:}
\begin{itemize}
\item \textbf{Complex calculations:} Conditional probabilities, combinatorics
\item \textbf{Simulation:} Generate random data, bootstrap samples
\item \textbf{Visualization:} Create probability distributions
\item \textbf{Problem solving:} Step-by-step solutions
\item \textbf{Code generation:} R scripts for probability problems
\end{itemize}

\vspace{0.5cm}
\textbf{Example prompts:}
\begin{itemize}
\item "Calculate the positive predictive value for a test with 95\% sensitivity, 90\% specificity, and 5\% prevalence"
\item "Simulate 10,000 binomial trials with n=20, p=0.3 in R"
\item "Explain the difference between relative risk and odds ratio with a medical example"
\end{itemize}
\end{frame}

\begin{frame}[fragile]
\frametitle{AI Demo: Diagnostic Test Calculations}
\textbf{Prompt:} "Create R code to calculate diagnostic test metrics"

\footnotesize
\begin{verbatim}
# Diagnostic Test Calculator
diagnostic_metrics <- function(sens, spec, prev) {
  # Calculate PPV and NPV
  ppv <- (sens * prev) /
         (sens * prev + (1-spec) * (1-prev))
  npv <- (spec * (1-prev)) /
         ((1-sens) * prev + spec * (1-prev))

  # Display results
  cat("PPV:", round(ppv*100, 1), "%\n")
  cat("NPV:", round(npv*100, 1), "%\n")

  return(list(PPV=ppv, NPV=npv))
}

# Example: COVID test
diagnostic_metrics(sens=0.85, spec=0.97, prev=0.05)
\end{verbatim}
\end{frame}

\begin{frame}[fragile]
\frametitle{AI Demo: Probability Simulations}
\textbf{Prompt:} "Simulate coin flips to estimate probability"

\footnotesize
\begin{verbatim}
# Simulate binomial trials
simulate_treatment <- function(n_patients=100,
                              p_success=0.7) {
  outcomes <- rbinom(n_patients, 1, p_success)
  success_rate <- mean(outcomes)

  cat("Simulated success rate:",
      round(success_rate*100, 1), "%\n")
  cat("Expected success rate:", p_success*100, "%\n")

  # Create visualization
  barplot(table(outcomes),
          names.arg=c("Failure", "Success"),
          col=c("red", "green"),
          main="Treatment Outcomes")
}

simulate_treatment()
\end{verbatim}
\end{frame}

\begin{frame}[fragile]
\frametitle{AI Demo: Visualizing Distributions}
\textbf{Prompt:} "Compare binomial and normal distributions in R"

\footnotesize
\begin{verbatim}
# Compare Binomial to Normal
n <- 100; p <- 0.3
x <- 0:n

# Binomial probabilities
binom_probs <- dbinom(x, n, p)

# Normal approximation
mu <- n * p
sigma <- sqrt(n * p * (1-p))

# Plot comparison
plot(x, binom_probs, type="h", col="blue",
     main="Binomial vs Normal",
     xlab="x", ylab="Probability")
curve(dnorm(x, mu, sigma), add=TRUE, col="red")
legend("topright", c("Binomial", "Normal"),
       col=c("blue", "red"), lty=1)
\end{verbatim}
\end{frame}

\begin{frame}
\frametitle{AI Best Practices for Probability}
\begin{columns}
\begin{column}{0.5\textwidth}
\textbf{DO:}
\begin{itemize}
\item Verify calculations manually
\item Check assumptions are met
\item Request step-by-step solutions
\item Ask for multiple approaches
\item Validate with simulations
\item Cross-reference formulas
\end{itemize}
\end{column}

\begin{column}{0.5\textwidth}
\textbf{DON'T:}
\begin{itemize}
\item Accept complex calculations blindly
\item Skip understanding the logic
\item Ignore edge cases
\item Misinterpret conditional probability
\item Confuse probability types
\item Trust without verification
\end{itemize}
\end{column}
\end{columns}

\vspace{0.5cm}
\textbf{Remember:} AI can calculate quickly, but you need to understand the concepts!
\end{frame}

\begin{frame}[fragile]
\frametitle{AI Assignment: Probability Problem Solving}
\textbf{This week's AI-integrated homework:}
\begin{enumerate}
\item Use AI to solve a diagnostic test problem, then verify manually
\item Generate R code for simulating coin flips, then modify for medical scenario
\item Ask AI to explain conditional probability three different ways
\item Create visualizations of different probability distributions
\item Identify and correct an AI error in probability calculation
\end{enumerate}

\vspace{0.5cm}
\textbf{Sample Problem:}
\small
\begin{verbatim}
"A new diagnostic test for Disease X has 92% sensitivity
and 88% specificity. In a population where 3% have the
disease, what is the probability that someone with a
positive test actually has the disease? Create R code
to calculate this and visualize how the answer changes
as prevalence varies from 1% to 10%."
\end{verbatim}
\end{frame}

\defaultcolors

\section{R Examples and Exercises}

\begin{frame}[fragile]
\frametitle{R Example: Basic Probability}
\small
\begin{verbatim}
# Probability calculations in R
# Sample space for rolling two dice
dice1 <- 1:6
dice2 <- 1:6
sample_space <- expand.grid(dice1, dice2)
sample_space$sum <- sample_space$Var1 + sample_space$Var2

# Probability of sum = 7
p_seven <- sum(sample_space$sum == 7) / nrow(sample_space)
cat("P(sum = 7) =", p_seven, "\n")  # 6/36 = 1/6

# Conditional probability example
# P(sum > 8 | first die = 4)
subset_data <- sample_space[sample_space$Var1 == 4, ]
p_conditional <- sum(subset_data$sum > 8) / nrow(subset_data)
cat("P(sum > 8 | first = 4) =", p_conditional, "\n")  # 3/6 = 1/2
\end{verbatim}
\end{frame}

\begin{frame}[fragile]
\frametitle{R Example: Medical Testing}
\footnotesize
\begin{verbatim}
# Breast cancer screening example
sens <- 0.90  # Sensitivity
spec <- 0.95  # Specificity
prev <- 0.008 # Prevalence

# Calculate PPV and NPV
ppv <- (sens * prev) /
       (sens * prev + (1-spec) * (1-prev))
npv <- (spec * (1-prev)) /
       ((1-sens) * prev + spec * (1-prev))

cat("PPV:", round(ppv * 100, 1), "%\n")
cat("NPV:", round(npv * 100, 1), "%\n")

# Simulate population
n <- 10000
n_disease <- round(n * prev)
n_healthy <- n - n_disease

# Test outcomes
TP <- round(n_disease * sens)
FP <- round(n_healthy * (1-spec))

cat("\nIn 10,000 women screened:\n")
cat("True positives:", TP, "\n")
cat("False positives:", FP, "\n")
\end{verbatim}
\end{frame}

\begin{frame}[fragile]
\frametitle{R Example: Binomial Distribution}
\small
\begin{verbatim}
# Clinical trial with 20 patients, treatment success rate 70%
n <- 20
p <- 0.7

# Probability of exactly 15 successes
p_15 <- dbinom(15, n, p)
cat("P(X = 15) =", round(p_15, 4), "\n")

# Probability of at least 15 successes
p_atleast_15 <- pbinom(14, n, p, lower.tail = FALSE)
cat("P(X >= 15) =", round(p_atleast_15, 4), "\n")

# Expected value and standard deviation
expected <- n * p
std_dev <- sqrt(n * p * (1-p))
cat("Expected successes:", expected, "\n")
cat("Standard deviation:", round(std_dev, 2), "\n")

# Visualize distribution
x <- 0:n
probs <- dbinom(x, n, p)
barplot(probs, names.arg = x, col = "lightblue",
        main = "Binomial Distribution (n=20, p=0.7)",
        xlab = "Number of Successes", ylab = "Probability")
abline(v = expected + 0.5, col = "red", lwd = 2, lty = 2)
\end{verbatim}
\end{frame}

\begin{frame}[fragile]
\frametitle{R Example: Normal Distribution}
\footnotesize
\begin{verbatim}
# Blood pressure in adults
mu <- 120      # Mean SBP
sigma <- 15    # Standard deviation

# Proportion with hypertension (SBP > 140)
p_hyper <- pnorm(140, mu, sigma, lower.tail = FALSE)
cat("P(BP > 140):", round(p_hyper * 100, 1), "%\n")

# Find percentiles
bp_90 <- qnorm(0.90, mu, sigma)
bp_95 <- qnorm(0.95, mu, sigma)
cat("90th percentile:", round(bp_90, 1), "mmHg\n")
cat("95th percentile:", round(bp_95, 1), "mmHg\n")

# Visualize distribution
x <- seq(70, 170, length = 100)
plot(x, dnorm(x, mu, sigma), type = "l",
     main = "Systolic Blood Pressure",
     xlab = "SBP (mmHg)", ylab = "Density")
abline(v = c(mu, 140), col = c("blue", "red"))
\end{verbatim}
\end{frame}

\begin{frame}
\frametitle{Practice Problems}
\begin{enumerate}
\item \textbf{Diagnostic Test:} A rapid strep test has 85\% sensitivity and 95\% specificity. If 15\% of patients with sore throat have strep, what's the NPV?

\item \textbf{Risk Calculation:} In a study of 1000 people, 200 smokers and 800 non-smokers, 30 smokers and 20 non-smokers developed heart disease. Calculate the relative risk.

\item \textbf{Binomial:} If a treatment works in 60\% of patients, what's the probability that it works in at least 8 out of 10 patients?

\item \textbf{Normal:} If cholesterol levels are normally distributed with mean 200 mg/dL and SD 40 mg/dL, what proportion of people have levels between 180 and 220?
\end{enumerate}
\end{frame}

\begin{frame}
\frametitle{Summary: Chapter 6 - Probability}
\textbf{Today we covered:}
\begin{itemize}
\item \textbf{Section 6.1:} Basic probability concepts, rules, and definitions
\item \textbf{Section 6.2:} Conditional probability and independence
\item \textbf{Section 6.3:} Diagnostic test applications (sensitivity, specificity)
\item \textbf{Section 6.4:} Random variables and expected value
\item \textbf{Section 6.5:} Discrete distributions (binomial)
\end{itemize}

\vspace{0.5cm}
\textbf{Key Learning Objectives:}
\begin{itemize}
\item Calculate and interpret basic probabilities
\item Understand conditional probability and independence
\item Apply probability concepts to diagnostic testing
\item Use probability distributions for modeling
\end{itemize}

\vspace{0.5cm}
\textbf{Next Week:} Chapter 7 \& 8 - Theoretical distributions and sampling
\end{frame}

\begin{frame}
\frametitle{Questions?}
\begin{center}
\Large{Thank you!}

\vspace{1cm}

\textbf{Office Hours:} Wednesday, 10:00 am -- 11:30 am\\
\textbf{Email:} john.molitor@oregonstate.edu\\
\textbf{Office:} 153 Milam

\vspace{1cm}

\textbf{Reading for next week:}\\
Pagano \& Gauvreau, Chapter 7 (7.1, 7.2, 7.4) and Chapter 8
\end{center}
\end{frame}

\end{document}